\section{BabyLM: cuestionar el culto a la escala mediante restricciones de datos}

La tercera unidad devuelve la atención de inference al preentrenamiento. BabyLM Challenge no busca el modelo más grande, sino que se limita a un presupuesto de datos cercano a la exposición lingüística de un niño y compara qué architecture, objective y curriculum aprenden con mayor eficiencia. Es a la vez un sample-efficiency benchmark y una reflexión sobre la concentración de los recursos de investigación y sobre los mecanismos del aprendizaje del lenguaje en los seres humanos.

\subsection{La restricción central del Challenge}

La diapositiva de título marca el paso de los reasoning methods a la data efficiency. El contexto de la clase es que el costo de entrenar modelos grandes sigue en aumento y que a las personas y a los laboratorios universitarios les resulta difícil participar en el frontier pretraining; por eso, fijar un presupuesto de datos pequeño vuelve a abrir un espacio de investigación comparable.

\lecturefigure{slide-24-babylm-challenge.jpg}{BabyLM Challenge estudia el aprendizaje eficiente del lenguaje con un presupuesto de datos pequeño}{Video de la clase de Stanford 00:15:15}

\begin{knowledgebox}{Lectura de la figura: la restricción misma cambia la pregunta de investigación}
Un scaling benchmark común pregunta «¿hasta dónde se llega con más data/params?»; BabyLM pregunta «¿cómo se aprovechan mejor los datos bajo un data cap?». Esto premia el inductive bias, el objective, la tokenización, el curriculum y la architecture, y no solo la obtención de más páginas web y más cómputo.
\end{knowledgebox}

\subsection{Qué es un developmentally plausible data budget}

La clase describe BabyLM así: entrenar modelos de lenguaje más pequeños con una cantidad de datos lingüísticos cercana a la que un niño puede recibir durante su crecimiento. La slide contrasta los billones de palabras del nivel de Chinchilla con los varios millones de palabras al año que recibe un niño, y el objetivo es conservar las capacidades gramaticales, semánticas y de razonamiento con muchos menos datos. Esta sección explica primero esa restricción de orden de magnitud y después por qué solo es developmentally inspired y no un modelo completo del aprendizaje infantil.

\lecturefigure{slide-25-what-is-babylm.jpg}{BabyLM contrasta el volumen de corpus de los modelos grandes con la exposición lingüística de un niño}{Video de la clase de Stanford 00:15:48}

\begin{knowledgebox}{Lectura de la figura: la comparación de volumen de datos es una referencia de orden de magnitud}
Las estimaciones de exposición léxica infantil dependen de la edad, de la familia y del método estadístico; además, un token de un LLM no equivale a una word que oye una persona. Esta comparación sirve para subrayar que la eficiencia de muestreo del preentrenamiento moderno es muy baja; no debe usarse para afirmar que un modelo de texto y un niño tienen la misma percepción, interacción u objetivos de aprendizaje.
\end{knowledgebox}

\begin{importantbox}{La eficiencia de muestreo debe reportarse junto con el presupuesto}
Si el desempeño promedio del modelo en el conjunto de tareas es $S$ y el número de tokens de entrenamiento es $N$, puede usarse
\[
\eta=\frac{S-S_0}{\log(1+N)}
\]
como una perspectiva aproximada de efficiency, donde $S_0$ es la línea base aleatoria o de un modelo pequeño. $\eta$ no es una métrica oficial de BabyLM; solo indica que al comparar deben reportarse a la vez la performance y el data budget, en lugar de mirar únicamente la puntuación absoluta.
\end{importantbox}

\subsection{Por qué vale la pena hacer BabyLM}

La slide Why BabyLM enumera la eficiencia, la apertura de nuevos casos de uso, la interpretabilidad, el alignment, el university access, la ciencia cognitiva y el aprendizaje humano del lenguaje. Convierte la investigación con modelos pequeños, de «sustituto para quien no puede costear modelos grandes», en una pregunta científica propia: qué datos y qué estructuras producen una generalization eficiente.

\lecturefigure{slide-26-why-babylm.jpg}{Motivaciones de BabyLM: eficiencia, accesibilidad e interdisciplina}{Video de la clase de Stanford 00:17:00}

\begin{knowledgebox}{Lectura de la figura: las siete razones se dividen en dos grupos, ingeniería y ciencia}
Del lado de la ingeniería importan el costo de entrenamiento, el despliegue y los presupuestos universitarios; del lado científico, la interpretabilidad, el alignment, la cognitive science y la child language acquisition. Con modelos pequeños es más fácil hacer una controlled ablation, pero tener pocos parámetros no implica automáticamente que el modelo sea interpretable o seguro.
\end{knowledgebox}

\teachervoice{El ponente enfatiza que la investigación con modelos grandes depende cada vez más de instituciones que poseen cientos o miles de GPU y millones de dólares. Uno de los valores de BabyLM es que las universidades y las personas vuelvan a poder entrenar modelos completos y comparar hipótesis, en lugar de solo invocar API cerradas.}

\subsection{Composición de los datos y límites de la analogía infantil}

La slide final sobre datos muestra la developmentally plausible pretraining mixture, que incluye fuentes como CHILDES child-directed speech, OpenSubtitles, libros infantiles, BNC speech, Wikipedia y Simple Wikipedia. Las distintas fuentes aportan diálogo, relatos, gramática cotidiana y conocimiento enciclopédico, pero siguen siendo una colección de textos fuera de línea.

\lecturefigure{slide-27-babylm-data.jpg}{La developmentally inspired pretraining data mixture de BabyLM}{Video de la clase de Stanford 00:19:09}

\begin{knowledgebox}{Lectura de la figura: la mixture simula tipos de entrada, no un entorno de desarrollo completo}
Child-directed speech aporta un estilo lingüístico interactivo, los subtitles aportan diálogo, los books aportan narrativa y Wikipedia aporta hechos y oraciones largas. Un niño real cuenta además con visión, acción, atención conjunta, corrección y retroalimentación social; una mixture de texto solo puede restringir el volumen de entrada lingüística, no reproducir el embodied learning.
\end{knowledgebox}

\begin{warningbox}{Developmentally plausible no equivale a cognitively equivalent}
El objective de entrenamiento, las epochs y el orden con que el modelo recibe el corpus completo difieren del aprendizaje continuo de un niño; al modelo también le faltan un cuerpo y objetivos sociales. BabyLM permite poner a prueba el sample-efficient language modeling, pero por sí solo no puede demostrar teorías de la adquisición humana del lenguaje.
\end{warningbox}

\teachervoice{El ponente espera un intercambio en ambos sentidos entre la ciencia cognitiva y el NLP: el aprendizaje eficiente de los humanos puede inspirar modelos, y los experimentos controlados con modelos pequeños también pueden, a su vez, poner a prueba hipótesis sobre el aprendizaje del lenguaje. El valor está en contar con una plataforma de experimentos comparables, no en reducir al niño a un next-token predictor.}

\subsection{Resumen de la sección}

BabyLM estudia la sample efficiency con presupuestos del orden de sub-100M-word, impulsa la comparación de architecture, objective, curriculum y composition de los datos, y reduce la barrera de entrada a la investigación. La analogía infantil ofrece una referencia de volumen y de tipos de datos, pero no incluye los mecanismos completos de percepción, interacción y desarrollo.
